\documentclass[10pt]{article}
\author{Yi Gu}


\usepackage{mathrsfs}

%\usepackage{amssymb,amsthm,amsmath,amssymb,latexsym,amsfonts,txfonts, pxfonts,wasysym,stmaryrd}
\textheight=230mm \textwidth=140mm
\topmargin=-1.5cm
\oddsidemargin=1.2cm
\evensidemargin=1.2cm
\setlength\parindent{0pt}
\parskip=7pt
%\setcounter{page}{1}
%\setlength{\baselineskip}{24pt} \baselineskip=24pt
\usepackage{amsmath,amsthm,amssymb}
\usepackage{graphicx}
\usepackage{bbm}
\usepackage{xcolor}
\usepackage[utf8]{inputenc}
\usepackage[english]{babel}
\usepackage{enumitem}
\usepackage{hyperref}
\usepackage{color}
\usepackage{colortbl}
\usepackage{tikz}
%\usepackage{vector}

\hypersetup{
    colorlinks=true,
    linkcolor=blue,
    filecolor=magenta,      
    urlcolor=cyan,
    pdftitle={Sharelatex Example},
    citecolor=red,
    % bookmarks=true,
}


\font\tencmmib=cmmib10 \skewchar\tencmmib '60
\newfam\cmmibfam
\textfont\cmmibfam=\tencmmib
\def\Ph{\mathchar'4010}
\def\Ps{\mathchar'4011}
\font\tensmc=cmcsc10
\def\smc{\tensmc}
\def\srm{\sevenrm }
\font\tenmsb=msbm10
\font\eightmsb=msbm10 scaled 1100
\def\Bbb#1{\hbox{\tenmsb#1}}
\def\Bbbb#1{\hbox{\eightmsb#1}}
\def\bbox{\quad\hbox{\vrule \vbox{\hrule \vskip2pt \hbox{\hskip2pt
\vbox{\hsize=1pt}\hskip2pt} \vskip2pt\hrule}\vrule}}
\def\lessim{\ \lower4pt\hbox{$
\buildrel{\displaystyle <}\over\sim$}\ }
\def\gessim{\ \lower4pt\hbox{$\buildrel{\displaystyle >}
\over\sim$}\ }
\def\n{\noindent}
\def\P{{\cal P}}
\def\si{\bar{\sigma}}
\def\E{{\cal E}}
\def\LL{{\cal L}}
\def\FF{{\cal F}}
\def\SS{{\cal S}}
\def\C{{\cal C}}
\def\D{{\cal D}}
\def\H{{\cal H}}
\def\X{{\cal X}}
\def\Y{{\cal Y}}
\def\A{{\cal A}}
\def\B{{\cal B}}
\def\O{{\cal O}}
\def\M{{\cal M}}
\def\I{{\cal I}}
\def\eps{{\varepsilon}}
\def\ch{{\mbox{\rm ch}\hspace{0.4mm}}}
\def\sh{{\mbox{sh}}}
\def\th{{\mbox{th}}}
\def\Av{{\mbox{\rm{Av}}}}
\def\n{\noindent}
\def\bbe{\vskip .15pc\hangindent=.5pc\hangafter=1}
\def\goin{\to\infty}
\def\qed{\hfill\break\rightline{$\bbox$}}



\newcommand{\e}{\mathbb{E}}
\newcommand{\p}{\mathbb{P}}
\newcommand{\lra}{\leftrightarrow}
\newcommand{\nlra}{\nleftrightarrow}
\newcommand{\bet}{\begin{theorem}}
\newcommand{\ent}{\end{theorem}}
\newcommand{\Var}{\mathrm{Var}}
\newcommand{\g}{\Gamma}
\newcommand{\de}{\delta}
\newcommand{\norm}[1]{\left\lVert#1\right\rVert}
\newcommand{\gy}[1]{\textbf{\color{red}[Yi: #1]}}
\newcommand{\olsi}[1]{\,\overline{\!{#1}}}
\newcommand\y{\cellcolor{blue!20}}
\newcommand\x{\cellcolor{red!20}}
\DeclareMathOperator{\Tr}{Tr}

\newtheorem{lemma}{\bf Lemma}
\newtheorem{claim}{\bf Claim}
\newtheorem{definition}{\bf Definition}
\newtheorem{theorem}{\bf Theorem}
\newtheorem{corollary}{\bf Corollary}
\newtheorem{remark}{\bf Remark}
\newtheorem{example}{\bf Example}
\newtheorem{exercise}{\bf Exercise}
\newtheorem{proposition}{\bf Proposition}
\newtheorem{question}{\bf Question}
\newtheorem{acknowledgements}{\bf Acknowledgements}

\newenvironment{Proof of lemma}{\noindent{\bf Proof of Lemma}}{\hfill$\Box$\newline}
\newenvironment{Proof of claim}{\noindent{\bf Proof of claim}}{\hfill$\Box$\newline}
\newenvironment{Proof of theorem}{\noindent{\bf Proof of Theorem}}{\hfill\scriptsize{$\Box$}\newline}
\newenvironment{Proof of theorems}{\noindent{\bf Proof of Theorems}}{\hfill$\Box$\newline}
\newenvironment{Proof of proposition}{\noindent{\bf Proof of Proposition}}{\hfill$\Box$\newline}
\newenvironment{Proof of propositions}{\noindent{\bf Proof of Propositions}}{\hfill$\Box$\newline}
\newenvironment{Proof of exercise}{\noindent{\it Proof of Exercise:}}{\hfill$\Box$}
\newenvironment{Acknowledgements}{\noindent{\bf Acknowledgements}}



\begin{document}
Hello Tuca, this is a short supplement document regarding our attempt on the asymptotic analysis of the following integral.
\begin{align} \label{target}
    \frac{1}{n}\log \int_{[0,\infty)^n} \exp \left(-x^T M_n x/2\right)dx.
\end{align}
Here $n = \binom{N}{2}$. The matrix $M_n$ has row and column indexed by 2-subsets of $[N]$. Moreover, $M_n$ satisfies the following properties:

1. For $N$ very large, $M_n$'s diagonal terms are asymptotically equal to $\frac{1}{4}$.

2. If $\{a,b\} \cap \{c,d\} = 1$, the entry at row $\{a,b\}$ and column $\{c,d\}$ is equal to $-\frac{1}{4N}$.

3. If $\{a,b\} \cap \{c,d\} = 0$, the entry at row $\{a,b\}$ and column $\{c,d\}$ is equal to $\frac{1}{2N^2}$.

(\ref{target}) can thus be rewritten as:
\begin{align}
    \binom{N}{2}^{-1}\log \int_{[0,\infty)^{\binom{N}{2}}} \exp \left(-x^T M_n x/2\right)dx.
\end{align}
One remark will be that the number of indices with empty intersection(the third property of $M_n$) is roughly $N$ times the number of those corresponds to the second property. However, the coefficients in front of entries in the second property are $O(N)$ times larger. Therefore, we should expect both types of entries will have equal level of impact in the computation of the integral.

Now we write the integral as an expectation.
\begin{align} \label{target_comp}
    &\int_{[0,\infty)^n} \exp \left(-x^T M_n x/2\right)dx \nonumber \\
    =&\int_{[0,\infty)^n} \exp \left(-\frac{1}{2}x^T M'_n x\right) \exp\left(-\frac{\norm{x}_2^2}{8}\right)dx \nonumber \\
    =&2^{-n}\int_{[0,\infty)^n} \exp \left(-2y^T M'_n y\right) \exp\left(-\frac{\norm{y}_2^2}{2}\right)dy \nonumber \\
    =&2^{-n}\left(\frac{\pi}{2}\right)^{n/2} \e \left( \exp \left(-2G^T M'_n G\right) \right),
\end{align}
where
\begin{align}
    M'_n = M_n - \frac{1}{4} \text{Id}_n,
\end{align}
and $G$ is a vector consists of i.i.d. standard half Gaussians. Note that the positive coefficients in the expectation come from those with second properties, rather than the third like in the original paper. Next up, we want to directly compute the rate function of $-G^T M'_n G$. If we manage to do so and it has a rate function $\mu_*(v)$, we can use truncation to get results similar to that from Varadhan's lemma. Indeed, we write the expectation in (\ref{target_comp}) as (my normalization scale could be wrong)
\begin{align}
    &\e \left( \exp \left(-2G^T M'_n G\right) \right) = \e \left( \exp \left(n \cdot 2\frac{-G^T M'_n G}{n}\right) \right) \nonumber \\
    \equiv& \e \left( \exp \left(n \cdot f(-G^T M'_n G)\right) \right).
\end{align}
We may expect a result like
\begin{align}
    \lim_{n \to \infty} \frac{1}{n} \log \e \left( \exp \left(-2G^T M'_n G\right) \right) = \sup_{v}\{2v - \mu_*(v)\}.
\end{align}
\end{document}